\vspace{-2mm}
\section{Conclusion}
\vspace{-2mm}

We introduce a comprehensive multi-dimensional benchmark, named \data~, dedicated to the evaluation of LVLMs, with a particular focus on measuring hallucinations in generative tasks. Our benchmark categorizes hallucinations into three distinct types -- object, attribute, and relation -- offering a detailed understanding of model inaccuracies. Furthermore, our novel evaluation framework, referred to as \eval~, employs a two-stage approach that integrates an LLM, effectively addressing the complexities related to open vocabularies, semantic similarities, and the intricate assessment of attributes and relationships. This method significantly enhances the precision and depth of image captioning evaluations compared to previous methods. Our experimental findings highlight the persistent challenges in this field, demonstrating that even state-of-the-art models such as GPT-4V, are prone to a considerable degree of hallucination. This study emphasizes the imperative for continuous advancements in LVLM evaluation techniques and establishes a new benchmark for future endeavors aimed at reducing hallucination and bolstering the reliability of content generated by LVLMs.